% !TeX program = xelatex
% Вопрос 1. Источники текста: ../tex/1/chapter.tex и библиографические сведения ../tex/1/main.tex.
% Ответ составлен только из цитат предоставленного TeX; по PDF оригинала в этой работе не сверялся.
\documentclass[12pt,a4paper]{article}
\usepackage{fontspec}
\setmainfont{Times New Roman}
\usepackage[english,russian]{babel}
\usepackage{amsmath,amssymb}
\usepackage[margin=22mm]{geometry}
\usepackage[unicode,hidelinks]{hyperref}
\setlength{\emergencystretch}{3em}
\setlength{\parindent}{0pt}
\renewcommand{\thefootnote}{\arabic{footnote})}
\newenvironment{textquote}{\begin{quote}}{\end{quote}}
\newcommand{\source}[2]{\par\nobreak\smallskip{\small Н. Е. Леонтьев. \textit{Основы теории фильтрации}. 2-е изд., 2017. Глава I, §~#1, с.~#2.}\par\medskip}
\begin{document}
% BEGIN QUESTION HEADER
\section*{Вопрос 1}
\noindent Основные физические сведения о пористых средах. Описание фильтрации в рамках механики сплошных сред. Пористость. Скорость фильтрации.

\medskip
\noindent\textbf{Источник:} Н. Е. Леонтьев. \textit{Основы теории фильтрации}. 2-е изд. Москва: МАКС Пресс, 2017. Глава I, §§ 1--2, печатные страницы 7--13. Ответ составлен из дословных цитат перенесённого в TeX текста учебника.
\par\bigskip
% END QUESTION HEADER
% Цитата 1; печатные страницы: 7
\begin{textquote}
\textit{Фильтрацией} называется движение (просачивание) жидкостей и газов через пористые среды под действием каких-либо факторов\footnote{Термин «фильтрация» также употребляется в более узком смысле как название процесса очищения жидкости или газа от примесей при течении через фильтр. В англоязычной литературе при описании этих процессов употребляется термин \textit{filtration}, в то время как о фильтрации в широком смысле этого слова говорят обычно как о «течении жидкости через пористую среду» (\textit{flow of a fluid through a porous medium}).}.
\end{textquote}
\source{1}{7}

% Цитата 2; печатные страницы: 7
\begin{textquote}
Вне зависимости от вызывающих фильтрацию причин, которые могут быть совершенно различными --- например, перепад давлений на границе среды, действие внешних массовых сил (сила тяжести, силы инерции, электромагнитные силы) или капиллярных сил на межфазных границах, деформация пористой среды и др.,--- речь всегда идет о движении жидкости по сложной системе разветвленных каналов, сообщающихся между собой. Как правило, внутренняя поверхность порового пространства имеет случайную структуру и не бывает точно известна, поэтому прямое описание движения жидкости во всех подробностях оказывается невозможным.
\end{textquote}
\source{1}{7}

% Цитата 3; печатные страницы: 7--8
\begin{textquote}
К счастью, однако, в большинстве приложений характерные линейные размеры рассматриваемых фильтрационных задач
много больше характерного размера пор, а при практическом изучении фильтрационных явлений обычно интересуются осредненными характеристиками течения. Это позволяет при описании крупномасштабных явлений рассматривать пористый материал как сплошную среду, характеристики которой (плотность и давление жидкости в порах и др.) в каждой точке получаются в результате осреднения по некоторой окрестности, содержащей достаточно большое число пор\footnote{Этот переход вполне аналогичен переходу от молекулярно-кинетической теории к континуальному описанию в случае, когда характерный размер явления много больше длины свободного пробега отдельной молекулы, с той лишь разницей, что в случае фильтрационных явлений описание течений на микроуровне тоже может быть осуществлено в рамках механики сплошных сред.}.
\end{textquote}
\source{1}{7--8}

% Цитата 4; печатные страницы: 8--9
\begin{textquote}
Например, при просачивании воды через пористый пласт под непроницаемой бетонной плотиной (рис.~1) характерный размер пор $d$ (доли миллиметра) много меньше характерного размера пористого пласта $L$ (десятки метров), поэтому мы можем перейти от описания течения на микроуровне (течение внутри конкретной поры) к описанию на макроуровне, считая, что, например, среднее давление $p$ жидкости в порах (в некотором месте внутри пласта в какой-то момент времени) получается осреднением давления в поровом пространстве внутри достаточно малого объема, размеры $l$ которого будут, с одной стороны, много меньше $L$ и, с другой стороны, много больше $d$ (скажем, порядка десятков сантиметров):
\begin{equation}
d\ll l\ll L.
\tag{1.1}\label{q1:eq-1-1}
\end{equation}
При этом, естественно, рассматриваемая величина (в нашем примере давление) будет, вообще говоря, некоторой функцией от
координат и времени: $p=p(x,y,z,t)$.
\end{textquote}
\source{1}{8--9}

% Цитата 5; печатные страницы: 9
\begin{textquote}
Важной характеристикой пористой среды является \textit{пористость} (или \textit{порозность}) $m$, равная относительной объемной доле порового пространства в материале. Пористость определяет количество жидкости, которое может содержаться в некотором объеме пористой среды (если жидкость целиком заполняет внутрипоровое пространство). Если для образца однородного пористого материала объемом $V$ объем пор составляет $V_{\text{п}}$, то пористость этого образца $m=V_{\text{п}}/V$.

Пористость, очевидно, является безразмерной величиной, значения которой лежат в интервале $0<m<1$. Для ориентировки укажем, что пористость почв лежит в диапазоне $0{,}3$--$0{,}7$, речного песка --- $0{,}3$--$0{,}55$, нефтегазоносных пластов, встречающихся на практике, --- $0{,}1$--$0{,}2$.
\end{textquote}
\source{2}{9}

% Цитата 6; печатные страницы: 9
\begin{textquote}
У некоторых пористых материалов часть пор изолирована от остального связного порового пространства или заполнена другой средой (например, при фильтрации воздуха в почве часть пор занята водой, удерживаемой в поровом пространстве за счет капиллярных сил). В подобной ситуации эти поры не участвуют в накоплении и перемещении интересующей нас жидкости внутри пористой среды, и поэтому вместо введенной выше величины (ее называют \textit{полной} или \textit{общей пористостью}) удобно использовать \textit{активную (эффективную) пористость}, равную объемной доле связанных между собой пор, которые могут быть заполнены жидкостью извне. […]
\end{textquote}
\source{2}{9}

% Цитата 7; печатные страницы: 9
\begin{textquote}
Для неоднородных пористых сред, свойства которых могут меняться от точки к точке, пористость $m$ будет некоторой функцией пространственных координат. Если пористая среда может деформироваться (это происходит, например, при утрамбовке грунта или изменении давления в жидкости в порах), то величина пористости может меняться и со временем: $m=m(x,y,z,t)$.
\end{textquote}
\source{2}{9}

% Цитата 8; печатные страницы: 10
\begin{textquote}
Подчеркнем, что при континуальном описании величина пористости (так же как и значения всех остальных осредненных параметров) определяется с помощью осреднения по области с размером порядка $l$, удовлетворяющим \eqref{q1:eq-1-1}.
\end{textquote}
\source{2}{10}

% Цитата 9; печатные страницы: 10
\begin{textquote}
Если сделать мысленный поперечный разрез образца пористой среды, то в образовавшемся сечении (обозначим его площадь $S$) часть площади $S_{\text{п}}$ будет приходиться на поры, а оставшаяся часть $S-S_{\text{п}}$ --- на твердое вещество пористой среды\footnote{«Каркас» пористой среды часто называют \textit{пористым скелетом} или \textit{матрицей}. Отметим также, что в отношении фильтрующейся среды (жидкой или газообразной) иногда применяется термин \textit{флюид}.}. Отношение $n$ площади пор в сечении к общей площади сечения называется \textit{просветностью (поверхностной пористостью)}: $n=S_{\text{п}}/S$.
\end{textquote}
\source{2}{10}

% Цитата 10; печатные страницы: 10--11
\begin{textquote}
Величина просветности зависит от того, через какую точку и в каком направлении проводится разрез среды, однако, что замечательно, \textit{среднее значение просветности} для большого набора параллельных сечений какого-то фиксированного образца равно
пористости этого образца.

[…]

Доказанное свойство позволяет приближенно считать пористость и просветность равными величинами для реальных пористых сред со случайной структурой: $m\approx n$.
\end{textquote}
\source{2}{10--11}

% Цитата 11; печатные страницы: 11
\begin{textquote}
Обратим внимание, что упоминавшиеся выше геометрические по своей природе параметры (пористость, характерный размер пор, функция распределения пор по радиусам) в общем случае не характеризуют полностью способность пористой среды пропускать через себя жидкость. Например, проницаемую для жидкости пористую среду можно превратить в непроницаемую, если ввести внутри порового пространства некоторое количество тонких перегородок, которые практически не поменяют перечисленные средние характеристики, но существенно изменят проводимость среды.
\end{textquote}
\source{2}{11}

% Цитата 12; печатные страницы: 11--12
\begin{textquote}
Для характеристики массопереноса при фильтрации жидкости (или газа) вводится векторная величина, называемая \textit{скоростью фильтрации}. Рассмотрим сначала простейший случай, когда жидкость движется вдоль тонкой прямолинейной трубки, заполненной пористым материалом. Для такого течения скорость фильтрации определяется как вектор $\boldsymbol{u}$, направленный в сторону движения
жидкости, величина которого равна объемному расходу жидкости (объему жидкости, протекающей в единицу времени) в расчете на единицу площади \textit{полного поперечного сечения} трубки (включающего как поры, так и пористую среду).

Здесь важно подчеркнуть, что хотя скорость фильтрации имеет размерность скорости, она тем не менее не равна скорости движения отдельных частиц жидкости. Действительно, если $S$ --- площадь поперечного сечения трубки, $S_{\text{п}}$ --- часть площади этого сечения, приходящаяся на поры, то постоянство объемного расхода однородной несжимаемой жидкости можно записать в виде
\[
Q=|\boldsymbol{u}|S=vS_{\text{п}}=\mathrm{const},
\]
где $Q$ --- объемный расход жидкости через трубку, $v$ --- среднее значение проекции \textit{скорости частиц жидкости} на ось трубки, вычисленное по площади сечения, занимаемой порами. Отсюда получаем
\[
v=\frac{|\boldsymbol{u}|S}{S_{\text{п}}}=\frac{|\boldsymbol{u}|}{n}\approx\frac{|\boldsymbol{u}|}{m},
\]
т.е. средняя скорость частиц жидкости приблизительно в $1/m$ раз больше скорости фильтрации.
\end{textquote}
\source{2}{11--12}

% Цитата 13; печатные страницы: 12
\begin{textquote}
В общем случае неодномерного движения жидкости в пористой среде \textit{скорость фильтрации} определяется как вектор $\boldsymbol{u}$, проекция которого на некоторое направление равна объемному расходу жидкости через единичную площадку, перпендикулярную данному направлению. С помощью этого определения легко вычисляются объемный расход $Q$ и массовый расход $Q_{\text{м}}$ жидкости через произвольную поверхность $\Sigma$, проведенную внутри пористой среды:
\begin{equation}
Q=\int_{\Sigma}\boldsymbol{u}\cdot\boldsymbol{n}\,d\sigma,
\qquad Q_{\text{м}}=\int_{\Sigma}\rho\boldsymbol{u}\cdot\boldsymbol{n}\,d\sigma,
\tag{2.1}\label{q1:eq-2-1}
\end{equation}
где $\rho=\rho(x,y,z,t)$ --- плотность фильтрующейся жидкости, $\boldsymbol{n}$ --- вектор нормали к поверхности $\Sigma$.

Удобство использования скорости фильтрации при практических исследованиях состоит в том, что с ее помощью можно находить расход жидкости без использования информации о пористости, которая явно не входит в соотношения \eqref{q1:eq-2-1}.
\end{textquote}
\source{2}{12}

\end{document}
